\section{UART 收发状态机、FIFO 与有界轮询}

CPU 向 THR 写入一个字节后，UART 先把它放入发送 FIFO。发送移位寄存器空闲
时取出字节，依次发出起始位、八个数据位和停止位。THRE 表示保持寄存器或 FIFO
可以继续接收，TEMT 表示 FIFO 与移位寄存器都空。两个状态对应提交与物理完成
的不同阶段。

115200 波特下每比特约 \(8.68\,\mu s\)，8N1 一帧十比特约
\(86.8\,\mu s\)。连续发送 32 字节理论线路时间约 \(2.78\,ms\)，远慢于
处理器执行几十条指令。FIFO 和状态轮询正是为了跨越这种速度差。

\topic{接收器用过采样抵抗相位误差}
接收端在检测到下降沿后，以内部多倍波特率时钟在比特中心附近采样。发送与接收
时钟存在少量误差，采样点会逐位漂移；停止位重新提供帧边界。帧越长，允许的
累计频差越小。

除数错误、线路噪声或格式不一致会产生帧错误、校验错误或错误字节。启动日志
只使用发送路径，但完整驱动以后要读取 LSR 错误位，并定义错误字节是丢弃、
上报还是保留带状态交给终端层。

\topic{FIFO 改变服务频率而不改变字符语义}
没有 FIFO 时，每个字符都要求软件及时读写寄存器。FIFO 缓冲一组字节，允许
中断在达到触发水位后再发生，减少处理器切换次数。FIFO 清空位通常是自清除
命令位，写入后不会像普通配置位一样长期保持。

当前初始化写 \code{0xC7}，同时启用 FIFO、清空收发队列并选择触发级别。
每个位都应有具名常量；组合值只出现在集中定义中。若硬件不支持 FIFO，驱动
应通过标识或行为降级，而不是假定写入一定生效。

\topic{modem control 保留了真实串行线路历史}
DTR 与 RTS 最初用于数据终端和调制解调器握手，OUT2 在 PC 上常参与 UART
中断门控。QEMU 的字符后端可能不需要真实调制解调器，但兼容寄存器仍按
16550A 行为实现。初始化这些位使代码与真实 PC UART 协议一致。

驱动不能把“QEMU 当前忽略某位”当成省略硬件步骤的依据。模型宽松时仍按规范
编写，故障测试则覆盖模型能够表达的错误边界。

\topic{轮询预算应从协议与阶段共同确定}
纯循环次数受处理器速度影响，不对应稳定墙钟。最早固件尚无定时器，只能使用
保守的大计数提供有限性；PIT 或 APIC timer 建立后，应把设备规范给出的时间
转换为单调时钟期限。

预算过短会在慢宿主或真实硬件上误报，过长会拖慢故障诊断。测试外层的 QEMU
超时必须显著大于目标内轮询正常完成时间，保证优先观察固件自己的失败分支，
外层只处理死循环和模拟器失控。

\section{日志协议、独立 oracle 与端到端证据}

纯文本方便人工阅读，但自动测试需要稳定标记。前缀标识项目和模块，阶段词标识
已建立的不变量，换行结束一条完整记录。后续字段可以采用
\code{key=value}，避免依赖自然语言句子。

\begin{osgridlongtable}{L{0.28\textwidth}L{0.31\textwidth}L{0.28\textwidth}}
\textbf{标记} & \textbf{能够证明} & \textbf{尚不能证明}\\ \hline
\endfirsthead
\textbf{标记} & \textbf{能够证明} & \textbf{尚不能证明}\\ \hline
\endhead
\code{RESET} & 入口、栈、初始化和首条发送可用 & 完整字符串与后续设备状态\\ \hline
\code{SERIAL_READY} & 多字节发送路径完成 & 磁盘、模式切换和内核存在\\ \hline
未来 \code{STAGE1} & 磁盘载入阶段已校验并执行 & 长模式与 ELF 正确\\ \hline
未来 \code{KERNEL64} & 64 位内核入口获得控制 & 用户态和文件系统存在\\ \hline
\end{osgridlongtable}

标记越晚，依赖链越长。失败时最后一条完整标记就是控制流达到的下界，不应把
缺失的下一标记直接解释为某个唯一根因。

\topic{测试 oracle 必须独立于被测实现}
若 ROM 审计器复用链接脚本中的同一公式，公式写错时两者可能一起通过。独立
oracle 直接从平台常量推导文件偏移，解码机器字节，再与规范目标比较。不同
实现路径降低共因错误。

系统测试不调用固件内部函数，而是启动 QEMU 并读取串口边界；单元测试对纯
地址计算使用宿主实现；随机测试生成大量非法位移。每层 oracle 的输入来源和
实现技术不同，组合后证据更强。

\topic{故障注入要穿过真实生产分支}
把测试函数直接写成“返回超时”只能验证上层处理，不能证明轮询循环会耗尽。
当前失败镜像通过构建期注入改变设备就绪观察，使同一发送过程走真实超时分支。
注入点位于硬件观察边界，生产控制流保持一致。

未来磁盘测试可让镜像截断、状态返回 ERR 或数据校验失败；页表测试可移除
present 位；系统调用测试可提供跨页非法缓冲。每个注入只改变一个前提，使
日志能够定位对应契约。

\topic{性质测试从不变量生成输入}
复位位移的性质是“只有解码后目标等于 \code{0xF000} 才接受”。生成器遍历或
随机抽取其余 16 位值，审计器必须全部拒绝。随机种子写入失败日志，任何反例
都可原样重放。

随机数量不是覆盖证明。生成域、过滤条件和断言必须来自规格；边界值仍用固定
样例明确覆盖。装载 ELF 时，可以围绕长度加法、表项数量、重叠和对齐构造
性质，而不是随意生成无法通过文件头的噪声。

\topic{端到端证据仍需二进制静态证据补强}
串口成功说明某条动态路径工作，不说明复位向量之外的填充、失败分支和全部符号
正确。构建验收同时保存镜像尺寸、关键偏移、ELF 机器类型、未解析符号和 QEMU
日志。

\begin{keyidea}
静态审计证明“产物具备预期结构”，动态测试证明“某状态下实际执行了预期路径”，
故障注入证明“前提破坏后系统按约定失败”。三类证据缺一不可。
\end{keyidea}
